% =====================================================================
%  单行本化之际的追记（原书 p.375–384 / PDF 385–394）
%  afterword —— 作者 2008 年为单行本所加的订正与补足（§5、§6 相关）
% =====================================================================

\chapter*{单行本化之际的追记}
\addcontentsline{toc}{chapter}{单行本化之际的追记（原书 p.375–384）}
\markboth{单行本化之际的追记}{单行本化之际的追记}

% [原文 p.375]
本书出版（1999 年）后发现的错误、不足之处等，记于此处。除此以外，只提及本书直接叙述的猜想等如今已被证明之类的情形。先声明：并未试图叙述 1999 年以后辛几何发展的重要部分。

\section*{A. 引理 5.83 的订正}

本书的引理 5.83 其实并不正确。问题出在取了由相当大的等价关系 5.82 所得的商空间这一点上。Gromov [158] 中出现了 geometric cycle 这一概念，但没有 geometric chain（几何链）这一概念。几何链这一概念与同调论的定义及基础密切相关，这里所写的做法并不准确。

下面要叙述的问题，同样也出现在 §6.6 的、例如定理 6.215 之中。另一方面，§6.1–§6.3 中所用的“几何链”不会引起问题。因为那里只处理成为 cycle 的情形。

一句话，若用 pseudo manifold 引入几何 cycle 这一概念，并把其间的“同调”仿照配边（但允许余维 2 的奇点）来定义，则把同调的几何 cycle 同一视之所得的商空间实现出同调群。这样虽能把握同调群，但要用“几何链”构成链复形更为困难，单纯用定义 5.82 的等价关系 5.82 作商是不行的。

% [Tutor Note: 原书 p.375 两处都印作“定義 5.82 の同値関係 5.82”，定义号与等价关系号被印成同一号码；此处照录，读者按“定义 5.82 中的那个等价关系”理解。]

幸而，从引理 5.83 到引理 5.98 的论证经若干修正便可订正。§6.6 中例如定理 6.215 的修正也同样可行。其概要说明如下。要点在于不用几何链，而用奇异链复形本身。此外，以下叙述依据 [122] 的修订版（[FOOO]）。

% [原文 p.376]
先作如下假定。

\begin{assumption}{A.1}
  若 $i>j$，则 $\mathcal{M}(R_i,R_j)$ 为空。
\end{assumption}

这只要按临界值 $f(R_i)$ 重新编号，便总能实现。

% [Tutor Note: 条件 A.1 的作用是让 $\mathcal{M}(R_i,R_j)$ 只在 $i<j$ 时非空，从而下文对 $i$ 的归纳法有意义（第 $i$ 步只需用到 $i$ 以下的项）。这不是几何上的限制，只是按临界值排序。]

$\mathcal{M}(R_i,R_j)$ 的紧化 $CM(R_i,R_j)$ 是带角流形。（参见引理 5.88。引理 5.88 的证明没有用到引理 5.83。）用 $S_mCM(R_i,R_j)$ 表示 $CM(R_i,R_j)$ 的余维 $m$ 角上的点全体。引理 5.88 的证明中，紧化 $CM(R_i,R_j)$ 的点被对各种 $k,l$ 的 $\mathcal{M}(R_k,R_l)$ 的纤维积层化（stratify），而 $m+1$ 个以上纤维积上的点全体就是 $S_mCM(R_i,R_j)$。这里证明下面的引理。（此引理与引理 5.84 的不同之处，在于不用几何链而用奇异链。）

\begin{lemma*}{A.2}
  存在由 $R_i$ 的光滑奇异链（smooth singular chain）组成的可数集合 $\mathcal{X}(R_i)$，满足下述条件。用 $C(R_i;\Z)$ 表示 $\mathcal{X}(R_i)$ 在奇异链复形 $S(R_i;\Z)$ 中生成的部分复形。
  \begin{enumerate}[label=（\arabic*）,leftmargin=2.4em]
    \item 从 $C(R_i;\Z)$ 到 $S(R_i;\Z)$ 的自然嵌入导出同调群的同构。
    \item 对 $\mathcal{X}(R_i)$ 的任意元素 $(P,\psi)$（这里 $P$ 是单形，$\psi\colon P\to R_i$ 是 $C^\infty$ 级映射），$\psi\colon P\to R_i$ 与对应于任意 $j$、$m=0,1,2,\dots$ 的
    \[
      \mathrm{ev}_{-,i,j}\colon S_mCM(R_i,R_j)\to R_i
    \]
    都横截。
    \item 对 $\mathcal{X}(R_i)$ 的任意元素 $(P,\psi)$，考虑纤维积
    \[
      P_\psi\times_{\mathrm{ev}_{-,i,j}}S_mCM(R_i,R_j).
    \]
    此时，选取流形 $P_\psi\times_{\mathrm{ev}_{-,i,j}}S_mCM(R_i,R_j)$ 的光滑三角剖分，可使每个单形与映射 $\mathrm{ev}_{+,i,j}$ 定出奇异单形。可以选取三角剖分，使这些奇异单形成为 $\mathcal{X}(R_j)$ 的元素。
  \end{enumerate}
\end{lemma*}

% [原文 p.377]
这里“$\psi\colon P\to R_i$ 与 $\mathrm{ev}_{-,i,j}$ 横截”是指：对 $P$ 的任何面 $P_a$（余维任意），限制 $\psi|_{P_a}\colon P_a\to R_i$ 都与 $\mathrm{ev}_{-,i,j}$ 横截。

\begin{proof}
  对 $i$ 作归纳来构成 $\mathcal{X}(R_i)$。满足 (1),(2) 的 $\mathcal{X}(R_1)$ 存在，是横截正则性定理的推论。

  作为归纳假设，假定下面的事。“对 $i\le i_0$ 的 $i$ 已选好 $\mathcal{X}(R_i)$，且满足 (1)。又对 $i\le i_0$ 的 $i$ 与任意 $j$，(2) 成立。进一步，对 $i,j\le i_0$ 的 $i,j$，(3) 成立。”

  在上述假定下构成 $\mathcal{X}(R_{i_0+1})$。

  首先，对 $(P,\psi)\in\mathcal{X}(R_i)$、$i\le i_0$，考虑纤维积
  \begin{equation}
    P_\psi\times_{\mathrm{ev}_{-,i,i_0+1}}S_mCM(R_i,R_{i_0+1}). \tag{A.1}
  \end{equation}
  由归纳假设与条件 A.1，此纤维积是横截的，因而定出流形。进而，映射
  \begin{equation}
    \mathrm{ev}_{+,i,i_0+1}\colon
    P_\psi\times_{\mathrm{ev}_{-,i,i_0+1}}S_mCM(R_i,R_{i_0+1})\to R_{i_0+1} \tag{A.2}
  \end{equation}
  对任意 $j$ 与 $m=0,1,2,\dots$，都与
  \begin{equation}
    \mathrm{ev}_{-,i_0+1,j}\colon S_mCM(R_{i_0+1},R_j)\to R_{i_0+1} \tag{A.3}
  \end{equation}
  横截。这由
  \[
    S_mCM(R_i,R_{i_0+1})\times_{R_{i_0+1}}S_{m'}CM(R_{i_0+1},R_j)
  \]
  是 $S_{m+m'+1}CM(R_i,R_j)$ 的某个成分，以及归纳假设得出。

  可以选取 (A.1) 的单形分割，使各单形上 $\mathrm{ev}_{+,i,i_0+1}$ 的限制与 (A.3) 横截。进一步，可以选取这样的单形分割，使它在 (A.1) 的边界上与已经选好的三角剖分一致。把 (A.1) 的每个单形经映射 $\mathrm{ev}_{+,i,i_0+1}$ 看作 $R_{i_0+1}$ 的奇异单形。把这样得到的奇异单形全体记作 $\mathcal{X}_0(R_{i_0+1})$。

  其次，按如下方式构成包含 $\mathcal{X}_0(R_{i_0+1})$ 的 $\mathcal{X}(R_{i_0+1})$。

  把满足定理 (2),(3) 的光滑奇异单形 $(P,\psi)$ 全体之集合记作 $\mathcal{X}_{\mathrm{trans}}(R_{i_0+1})$。其元素的 1 次结合构成奇异链复形的部分复形，记作 $S_{\mathrm{trans}}(R_{i_0+1};\Z)$。用横截正则性定理，可证明嵌入 $S_{\mathrm{trans}}(R_{i_0+1};\Z)\to S(R_{i_0+1};\Z)$ 在同调上导出同构。

  首先，把 $\mathcal{X}_1(R_{i_0+1})$ 取成

  % [原文 p.378]
  $\mathcal{X}_1(R_{i_0+1})\supset\mathcal{X}_0(R_{i_0+1})$、$\mathcal{X}_1(R_{i_0+1})\subset\mathcal{X}_{\mathrm{trans}}(R_{i_0+1})$，
  并且 $\mathcal{X}_1(R_{i_0+1})$ 生成部分复形 $C_1(R_{i_0+1})$，而嵌入导出的映射 $H(C_1(R_{i_0+1});\Z)\to H(R_{i_0+1};\Z)$ 是满射。

  $H(C_1(R_{i_0+1});\Z)\to H(R_{i_0+1};\Z)$ 未必是单射。

  于是把 $\mathcal{X}_2(R_{i_0+1})$ 取成
  \[
    \mathcal{X}_2(R_{i_0+1})\supset\mathcal{X}_1(R_{i_0+1}),\qquad
    \mathcal{X}_2(R_{i_0+1})\subset\mathcal{X}_{\mathrm{trans}}(R_{i_0+1}),
  \]
  并且 $\mathcal{X}_2(R_{i_0+1})$ 生成部分复形 $C_2(R_{i_0+1})$，而嵌入导出的映射 $H(C_1(R_{i_0+1});\Z)\to H(R_{i_0+1};\Z)$ 的核（kernel）被嵌入导出的映射 $H(R_{i_0+1};\Z)\to H(C_2(R_{i_0+1});\Z)$ 送到 $0$。

  若 $H(C_2(R_{i_0+1});\Z)\to H(R_{i_0+1};\Z)$ 的核不为 $0$，就同样地选取 $\mathcal{X}_3(R_{i_0+1})$。如此反复，
  \[
    \mathcal{X}(R_{i_0+1})=\bigcup_k\mathcal{X}_k(R_{i_0+1})
  \]
  便具有引理所断言的性质。
\end{proof}

\begin{remark}{A.3}
  前面说：取带角流形 (A.1) 的单形分割，用 (A.2) 便得到 $R_i$ 的奇异链。这多少不准确。即：对 (A.1) 中每个最高维单形，确定它与标准单形的同构，才能用 (A.2) 得到 $R_i$ 的奇异链。每个单形与标准单形的同构，一旦确定单形顶点的顺序便唯一确定。依次选取顶点的顺序、确定各单形与标准单形的同构，就定出 $R_i$ 的奇异链。顶点顺序的取法从边界起依次决定。即假定 (A.1) 边界上的顶点已有顺序，再把它扩充到内部。
\end{remark}

用引理所得的链复形 $C_*(R_i;\Z)$，令
\[
  S_0^k(R_i;\Z)=C_{\dim R_i-k}(R_i;\Z),
\]
并用它把定义 5.85 换成下面这样。

\begin{definition}{A.4}
  \[
    CF^k(X;f)=\bigoplus_i S_0^{k-\mu(u)}(R_i;\Z).
  \]
\end{definition}

% [原文 p.379]
边界算子 (5.17) 的定义对 $CF^k(X;f)$ 的链 $(P,\psi)$ 可以给出。这是引理 A.2 的 (2),(3) 的推论。由此构成满足 (5.19) 的谱序列。

要证明定理 5.90，仍然不用几何链而用奇异链。包括注 5.83 在内，论证几乎同样进行。不过引理 5.97 的证明需要小心。即：由像的维数比原本的维数（作为奇异链的维数）小 $1$，推不出作为奇异链为 $0$（在定义 5.82 中，作为 current 相同者视为相同，加入这个等价关系后就变成 $0$。但在我们现在的做法中不允许这样）。这里作如下处理。考虑引理 5.98 中用的 $\Phi=\sum_i\Phi_i$。它被用来作由 $\rho=0$ 时的 $f_{\tau,\rho}$ 得到的链映射
\begin{equation}
  (f_{\tau,0})_*\colon CF(X;f_{-\infty})\to CF(X;f_{-\infty}) \tag{A.4}
\end{equation}
与由 $\rho=1$ 时的 $f_{\tau,\rho}$ 得到的链映射
\begin{equation}
  (f_{\tau,1})_*\colon CF(X;f_{-\infty})\to CF(X;f_{-\infty}) \tag{A.5}
\end{equation}
之间的链同伦。引理 5.97 意味着 $(f_{\tau,0})_*$ 是恒等映射，而这就带有上述问题。

% [Tutor Note: (A.4),(A.5) 中的 $f_{-\infty}$、$f_{\tau,\rho}$ 是第 5 章 Floer 同调里哈密顿量的族；$(f_{\tau,\rho})_*$ 是它诱导的链映射。原证明想把 $\Phi$ 当作 $(f_{\tau,0})_*$ 到恒等的链同伦，而问题在于“像的维数少 1”在奇异链层面并不等于 $0$。]

于是把 $\Phi_i$ 稍作改动。即把定义中所用的空间
\begin{equation}
  \bigcup_\rho P_\psi\times_{\mathrm{ev}_{-,i}}
  \mathcal{M}(R_{i,-},R_{j,-};f_{\tau,\rho}) \tag{A.6}
\end{equation}
在 $i\ne j$ 时，只把 $\rho=0$ 处缩掉。即在彼处有 $\R$ 的作用，就用该作用缩掉。一旦知道缩掉后的空间具有三角剖分，由于 $\rho=0$ 处没有本应有的维数的单形，$i\ne j$、$\rho=0$ 所对应的部分作为奇异复形便成为 $0$。（这一点后面再更准确地叙述。）

以下把 $R_{i,-}$ 写作 $R_i$。观察方程（(5.23) 之后的常微分方程，p.227），可知 $i=j$ 时
\begin{equation}
  \mathcal{M}(R_i,R_i;f_{\tau,\rho}) \tag{A.7}
\end{equation}
当 $\rho\to0$ 时趋近 $R_i$ 本身，且 $\mathrm{ev}_{-,i,i},\mathrm{ev}_{+,i,i}$ 趋近恒等映射 $R_i\to R_i$（实际上，$i=j$ 时，(5.23) 之后的常微分方程的解在 $\rho$ 充分接近 $0$ 时处处位于 $R_i$ 附近。因此，

% [原文 p.380]
(A.7) 成为 $R_i$ 的稳定流形的小扰动与不稳定流形之交，即得上断言）。

这样便知 $\Phi$ 给出 (A.5) 与恒等映射之间的链同伦，定理 5.90 得证。

最后，就 (A.6) 用 $\R$ 作用缩掉这一构造再说几句。用 $\R$ 作用缩掉是看起来危险的构造；$\R$ 不紧，人们大概会担心它显然不再是 Hausdorff 空间。

为简单起见，考虑 $\mathcal{M}(R_i,R_j;f_{-\infty})$ 无边界的情形，即临界值 $f_{-\infty}(R_i)$ 与临界值 $f_{-\infty}(R_j)$ 之间没有 $f_{-\infty}$ 的临界值的情形。
\[
  \mathcal{M}(R_i,R_j;f_{\tau,0})
  =\tilde{\mathcal{M}}(R_i,R_j;f_{-\infty}),
\]
且
\begin{equation}
  \tilde{\mathcal{M}}(R_i,R_j;f_{-\infty})
  =\mathcal{M}(R_i,R_j;f_{-\infty})\times\R . \tag{A.8}
\end{equation}
叙述这个同构。令
\[
  f_{-\infty}(R_i)=c_i,\qquad f_{-\infty}(R_j)=c_j,
\]
取满足
\[
  c_i<c<c_j
\]
的 $c$。对 $\gamma\in\mathcal{M}(R_i,R_j;f_{\tau,0})$，用
\[
  f_{-\infty}(\gamma(r(\gamma)))=c
\]
定义 $r(\gamma)\in\R$。则同构 (A.8) 的第 2 分量由这个 $r$ 给出。

同样地，用
\[
  f_{\tau,\rho}(\gamma(r(\gamma)))=c
\]
定义
\[
  r\colon\mathcal{M}(R_i,R_j;f_{\tau,\rho})\to\R .
\]
对充分小的 $\rho$，$r$ 有定义。用它可定义微分同胚
\begin{equation}
  \mathcal{J}\colon\mathcal{M}(R_i,R_j;f_{-\infty})\times[0,\rho_0)\times\R
  \longrightarrow\bigcup_{\rho\in[0,\rho_0)}\mathcal{M}(R_i,R_j;f_{\tau,\rho}), \tag{A.9}
\end{equation}
第 2 分量 $\rho$ 在左右两边保持，且
\[
  r(\mathcal{J}(\gamma,\rho,s))=s .
\]
看第 3 分量（即 $s$）趋于 $\pm\infty$ 的极限时会发生什么。

% [原文 p.381]
$\rho>0$ 时
\begin{equation}
\begin{aligned}
  \partial\mathcal{M}(R_i,R_j;f_{\tau,\rho})
  ={}&\mathcal{M}(R_i,R_i;f_{\tau,\rho})_{\mathrm{ev}_{+,i}}
  \times_{\mathrm{ev}_{-,i,j}}\mathcal{M}(R_i,R_j;f_{-\infty})\\
  &\quad\cup\mathcal{M}(R_i,R_j;f_{-\infty})_{\mathrm{ev}_{+,i,j}}
  \times_{\mathrm{ev}_{-,j}}\mathcal{M}(R_j,R_j;f_{\tau,\rho})
\end{aligned}
\tag{A.10}
\end{equation}
成立。右边第 1 项、第 2 项分别对应 $s\to\infty$、$s\to-\infty$ 的极限。不过 (A.10) 在 $\rho=0$ 时没有意义。

但是，用 (A.10) 可得与 (A.9) 在公共部分一致的微分同胚
\[
  \mathcal{J}'\colon\mathcal{M}(R_i,R_j;f_{-\infty})\times(0,\rho_0)
  \times(\R\cup\{\pm\infty\})
  \longrightarrow\bigcup_{\rho\in(0,\rho_0)}CM(R_i,R_j;f_{\tau,\rho}).
\]
用 $\mathcal{J}$ 把 $\mathcal{M}(R_i,R_j;f_{\tau,0})$ 与 $\mathcal{M}(R_i,R_j;f_{-\infty})\times\R$ 同一视之，则 $\mathrm{ev}_{-,i},\mathrm{ev}_{+,j}$ 与第 2 分量 $s\in\R$ 无关。

现在考虑
\begin{equation}
  \bigcup_{\rho\in[0,\rho_0)}CM(R_i,R_j;f_{\tau,\rho}) \tag{A.11}
\end{equation}
中，$\rho=0$ 时把仅 $s$ 坐标不同的两点同一视之所得者。它可与
\begin{equation}
  \mathcal{M}(R_i,R_j;f_{-\infty})\times[0,\epsilon)^2 \tag{A.12}
\end{equation}
同一视之。(A.11) 的 $\rho=0$ 部分对应 (A.12) 第 2 分量为 $(0,0)$ 的部分。利用 $\rho=0$ 时 $\mathrm{ev}_{-,i},\mathrm{ev}_{+,j}$ 与第 2 分量 $s\in\R$ 无关这一点，可把 $\mathrm{ev}_{-,i},\mathrm{ev}_{+,j}$ 看作来自 (A.12) 的映射（要找出使该映射在顶点处可微的微分结构，需一些论证。可与 [FOOO] 的附录作同样的讨论）。

把 (A.12) 与 $(P,\psi)$ 取纤维积，便得到把 (A.6) 缩掉后具有所需性质的空间。

% [Tutor Note: 这一段的要点是：把 $\R$ 作用“缩掉”会让一部分单形失去本应有的维数，从而在奇异链层面自动为 $0$；但代价是必须在 $(0,\rho_0)$ 与 $\rho=0$ 的接缝处另找微分结构（即 (A.12) 的角），这正是 [FOOO] 附录处理的技术点。]

\section*{B. 更细微之处的订正补足}

1) 关于 Weinstein 猜想 4.25

% [原文 p.382]
这个猜想原先附加了假定 $H^1(M;\R)=0$，这并不必要。$M$ 为 3 维时，该猜想最近由 Taubes [Tau] 证明。

2) 关于定理 5.121

这个定理的假定原写作“半正”，但 [122] 的证明有一处缺口，在现行的修订版（[FOOO]）中，假定是这样的：$M$ 是满足下列之一的一些 $M_i$ 的直积。
\begin{enumerate}[label=（\arabic*）,leftmargin=2.4em]
  \item $\dim_\C M_i\le2$。
  \item $M_i$ 是单调的。
  \item $M_i$ 是 Fano 流形。
\end{enumerate}
不过，猜想 5.120 仍被认为正确。

3) 关于 $\C^n$ 的拉格朗日子流形

定理 5.126 如今已被改进，例如已知：若 $L$ 是自旋流形且 $\pi_k(L)=0$（$k\ge2$），则存在 $\beta\in\pi_1(L)$ 使 $\mu(\beta)=2$。（参见 [Fuk2]。）

4) 关于“定理 6.129”

基于该定理的方略 6.130 的详细证明现已完成，在 [FOOO] 第 4 章。

5) 关于特殊拉格朗日纤维丛

现在人们不相信（在复 3 维以上）存在（只允许余维 2 奇点的）特殊拉格朗日纤维丛。人们认为它在以下意义下“渐近地存在”。

取 $M_\tau$ 中当 $\tau\to0$ 时趋近巨大复结构极限的族，则在 $S(N)$ 的 $o(\tau)$ 邻域之外，存在特殊拉格朗日纤维丛。

拉格朗日纤维丛在各种情形下已被构造出来。

6) 关于通量猜想

猜想 6.201 已由小野薫 [Ono] 证明（在这种形式下）。

7) 关于相对自旋结构

定理 6.213 中写道：由 $\tilde w_2$ 与 $L$ 的定向决定伪全纯圆盘模空间的定向。这多少不准确。准确地说，可以定义 $L$ 的相对自旋结构这一概念，$L$ 的定向与 $L$ 的相对自旋结构决定伪全纯圆盘模空间的定向。

% [原文 p.383]
若 $\tilde w_2$ 存在，则 $L$ 具有相对自旋结构，但 $\tilde w_2$ 的取法与 $L$ 的相对自旋结构并非一一对应。

例如，$w_2(L)=0$ 时，对应于 $\tilde w_2=0$ 的 $L$ 的相对自旋结构，与 $L$ 的自旋结构一一对应。

详见 [FOOO]。

8) 关于定义 6.214

由于本追记 A 中所述的理由，这个定义有问题。即：用几何链构成的链复形去作 $A_\infty$ 代数，不是正确的做法（它不给出正确的同调群）。正确的做法是用奇异链复形的某个应有的可数部分复形，这在 [FOOO] 的 §30 中有详细叙述。

也可以使用 de Rham 上同调（[Fuk3]）。还有 Joyce 的仓西同调理论 [Joy] 这种做法。

9) 关于引理 6.225（循环对称性）

用 [FOOO] 的 §30 实现 $A_\infty$ 结构时，对称性会破坏，该引理 6.225 在 [FOOO] 中并未被证明。

人们认为，用 de Rham 上同调代替奇异同调可使这一证明成为可能（参见 [Fuk3] 等），但细节尚未写出。据 Joyce，用仓西同调便可实现循环对称性。

10) 关于 $A_\infty$ 范畴的构成

关于注 6.231 中提到的 $A_\infty$ 范畴的构成，笔者写了 [Fuk1]。[Fuk1] 中省略了许多证明，但把它们按 [FOOO] 的做法稍作改动即可全部证明。

11) 关于“伪全纯曲线”这一名称

笔者在本书（1999 年版）中把「概正則曲線」用作 pseudo-holomorphic curve 的译语。这是为了与「概複素構造」的「概」取得一致，但其他

% [原文 p.384]
许多人使用「擬正則曲線」这一译语。作为 pseudo 的译语，确实「擬」更自然，故笔者今后也想改用「擬正則曲線」。（数学辞典第 4 版也译作「擬正則曲線」。）

% [Tutor Note: 第 11 条讨论的是日文译语「概正則曲線」与「擬正則曲線」之争，与中译名无关；本译注本全书统一作“伪全纯曲线”，故正文保留日文原词以便读者理解作者的语气。]

最后附上追加文献。追加的文献限于本追记中提到的。1999 年以后虽有许多进展，但未试图把它们纳入。

{\small
\noindent [Fuk1]\ K. Fukaya, Floer homology and mirror symmetry. II. Minimal surfaces, geometric analysis and symplectic geometry (Baltimore, MD, 1999), \textit{Adv. Stud. Pure Math.} \textbf{34}, pp. 31–127, Math. Soc. Japan, Tokyo, 2002.\par\smallskip

\noindent [Fuk2]\ K. Fukaya, Application of Floer homology of Lagrangian submanifolds to symplectic topology. in \textit{Morse theoretic methods in nonlinear analysis and in symplectic topology.} NATO Sci. Ser. II, \textit{Math. Phys. Chem.} \textbf{217}, pp. 231–276, Springer, Dordrecht, 2006.\par\smallskip

\noindent [Fuk3]\ K. Fukaya, Differentiable operad, Kuranishi correspondence, and Foundation of topological field theories based on pseudo-holomorphic curve, preprint.\par\smallskip

\noindent [FOOO]\ K. Fukaya, Y.-G. Oh, K. Ono and H. Ohta, Lagrangian intersection Floer theory: anomaly and obstruction.（2006 年版及 2008 年版的一部分放在 \url{http://www.math.kyoto-u.ac.jp/~fukaya/} 上。原稿在现时点（2008 年 9 月）已完成。）\par\smallskip

\noindent [Joy]\ D. Joyce, Kuranishi homology and Kuranishi cohomology, preprint, arXiv:0707.3572.\par\smallskip

\noindent [Ono]\ K. Ono, Floer–Novikov cohomology and the flux conjecture, \textit{Geom. Funct. Anal.} \textbf{5}(2006), pp. 981–1020.\par\smallskip

\noindent [Tau]\ C. Taubes, The Seiberg–Witten equations and the Weinstein conjecture, \textit{Geom. Topol.} \textbf{11}(2007), pp. 2117–2202.\par\smallskip
}

% CHUNK afterword DONE
